\subsection{Fonction de r\'{e}compense $R(s,a)$ --- Bilan co\^{u}ts/b\'{e}n\'{e}fices}

\paragraph{Cadre \'{e}conomique}
Conform\'{e}ment \`{a} §4.3, tous les param\`{e}tres sont exprim\'{e}s en FCFA et
sour\c{c}\'{e}s dans les textes officiels camerounais
(voir \texttt{data/references/baremes\_dgi\_mint.md}).

\begin{table}[H]
\centering
\caption{Montants de r\'{e}f\'{e}rence utilis\'{e}s dans $R(s,a)$ --- sources §4.3}
\label{tab:c:baremes}
\begin{tabular}{lrl}
\toprule
\textbf{Param\`{e}tre} & \textbf{Montant (FCFA)} & \textbf{Source} \\
\midrule
\multicolumn{3}{l}{\textit{Gains}} \\
Amende l\'{e}g\`{e}re (infraction doc) & 25\,000 & Code Route n°96/07, Art.~137 \\
Amende falsification & 500\,000 & Code Route n°96/07, Art.~169 \\
Vignette moyenne (parc cam.) & 45\,000 & Loi Finances DGI 2023, Art.~207--209 \\
Majoration retard 30j & +25\,\% & CGI Cameroun, Art.~556 \\
Saisie v\'{e}hicule (frais jud.) & 200\,000 & Code Proc. P\'{e}nale, Art.~78--80 \\
\midrule
\multicolumn{3}{l}{\textit{Co\^{u}ts op\'{e}rationnels}} \\
A0 --- Laisser passer & 500 & Bar\`{e}me MINT 2023 (5\,sec agent) \\
A1 --- Contr\^{o}le standard & 5\,000 & 5\,min agent + flux \\
A2 --- Arr\^{e}t + saisie & 50\,000 & 2 agents 30\,min + logistique \\
A3 --- Signalement DGI & 2\,000 & 1\,min agent + frais postaux \\
A4 --- Transfert PJ & 75\,000 & 2 agents 30\,min + dossier PJ \\
\midrule
\multicolumn{3}{l}{\textit{Risque de contestation}} \\
Co\^{u}t total contestation & 450\,000 & Tribunal Admin., Art.~12 \\
\bottomrule
\end{tabular}
\end{table}

\paragraph{Formule g\'{e}n\'{e}rale}
\begin{equation}
R(s,a) = \mathbb{E}[\text{gain}] - \text{co\^{u}t\_op\'{e}rationnel}
          - \mathbb{E}[\text{risque\_contestation}]
\end{equation}
\begin{equation}
= P(\text{fraude r\'{e}elle} \mid s) \times G_a
  - C_a - \bigl(1 - P(\text{fraude} \mid s)\bigr) \times \mathrm{Rq}_a
\end{equation}

o\`{u} $G_a$ est le gain maximal de l'action $a$, $C_a$ son co\^{u}t
op\'{e}rationnel et $\mathrm{Rq}_a$ le risque de contestation.
$P(\text{fraude} \mid s)$ est estim\'{e} par profil d'\'{e}tat à partir de
la note interne MINT/DGI §1.2 (taux 7--15\,\%).

\paragraph{Exemples illustratifs (jury §5)}
\begin{table}[H]
\centering
\caption{Exemples de calcul $R(s,a)$ --- valeurs r\'{e}elles du pipeline}
\label{tab:c:exemples_r}
\begin{tabular}{llrl}
\toprule
\textbf{\'{E}tat $s$} & \textbf{Action $a$} & \textbf{$R(s,a)$~(FCFA)} & \textbf{D\'{e}tail calcul} \\
\midrule
Conforme·Haute·$\neg$Alerte & A2 Arr\^{e}t+Saisie
  & -442\,500\,FCFA
  & $0.05 \times 700\,000 - 50\,000 - 0.95 \times 450\,000$ \\
D\'{e}grad\'{e}·Haute·$\neg$Alerte & A0 Laisser passer
  & -500\,FCFA
  & $-500\,\text{FCFA}$ (co\^{u}t agent uniquement) \\
Expir\'{e}·Moy·$\neg$Alerte & A3 Signalement DGI
  & 37\,375\,FCFA
  & $0.70 \times 56\,250 - 2\,000$ \\
\bottomrule
\end{tabular}
\end{table}

\paragraph{Interpr\'{e}tation}
La matrice $R$ encode la connaissance op\'{e}rationnelle MINT/DGI :
l'action la plus profitable d\'{e}pend de l'\'{e}tat --- il n'existe pas
d'action universellement optimale.
C'est pr\'{e}cis\'{e}ment pourquoi le Module~C n\'{e}cessite un MDP pour calculer
la politique $\pi^*$ qui maximise la r\'{e}compense cumul\'{e}e sur l'ensemble des
\'{e}tats (Modules~D : it\'{e}ration de valeur et it\'{e}ration de politique).
